%% Addition to Section 2: Background and Related Work
%% New subsection: Transformer Mechanistic Interpretability

\subsection{Transformer Mechanistic Interpretability}
\label{sec:mech-interp}

Recent work on mechanistic interpretability has revealed that transformer attention
heads perform specialized functions: induction heads~\cite{olsson2022induction}
that implement in-context learning by pattern-matching prior sequences, name
movers~\cite{wang2023interpretability} that propagate information about specific
entities through the residual stream, and safety
heads~\cite{zou2023representation} that encode refusal-related
representations in late layers. Our framework provides a geometric unification:
each head type corresponds to a standpoint layer, and the composition of heads
implements parallel transport on a fiber bundle. The key difference from prior
work is that we do not merely \emph{describe} what heads do --- we prove that
their collective operation \emph{is} parallel transport, with curvature measuring
standpoint consistency.

The key insight from mechanistic interpretability --- that the residual stream
is the central communication channel of the transformer, and that attention heads
read from and write to specific subspaces of this
stream~\cite{elhage2021mathematical} --- is precisely the structure we formalize.
Each head's value matrix $\mathbf{V}_h$ defines a subspace of the residual stream;
the collection of these subspaces defines the fiber decomposition; and the
attention weights define the transport between events.

Sheaf neural networks~\cite{hansen2020toward, bodnar2022neural} apply
sheaf-theoretic consistency conditions to graph-structured data. Our transport
consistency conditions are structurally related: both frameworks penalize
disagreement between local representations propagated along edges. The key
difference is the object being made consistent: sheaf networks enforce
representational agreement across node features for a fixed graph, whereas our
self-connection enforces standpoint agreement across causally ordered events in
a DAG, with the additional structure that disagreement is attributed to a
specific standpoint layer via the block decomposition of $\mathbf{F}_{ijk}$.

The Attention Schema Theory~\cite{graziano2017attention} explicitly targets AI
but focuses on attention modelling rather than standpoint continuity. Our
framework subsumes this: attention modelling is the transport mechanism, and
standpoint continuity is the curvature constraint. The Global Workspace
Theory~\cite{baars1988cognitive, dehaene2011experimental} provides the
information-broadcasting analogue; our salience misfire mode (F6) corresponds
to a failure of this broadcast --- a local event fails to be transported to the
global standpoint, introducing curvature in the narrative component.

Representation engineering~\cite{zou2023representation} identifies linear
directions in activation space that encode specific behaviors (e.g., refusal,
truthfulness). Our framework provides the geometric structure underlying these
directions: they are basis vectors of the standpoint subspaces $W_k$, and the
``refusal direction'' is a generator of the moral subspace $W_{\mathrm{mor}}$.
The difference is that representation engineering identifies individual
directions, whereas we identify the full subspace structure and its transport
properties.
